В ходе работы был рассмотрен подход к обработке видео на основе скриптового
фреймсервера AviSynth в связке с VirtualDub2. Модель фреймсервера, при
которой скрипт отдаёт кадры по требованию внешней программы, определяет
главное преимущество связки перед графическими видеоредакторами:
воспроизводимость~\cite{avisynth}. Один раз написанный текстовый скрипт
гарантирует, что цепочка фильтров будет применена идентично на любой
машине, а отсутствие промежуточных файлов сокращает объём дисковых
операций.

Компактность синтаксиса AviSynth делает его удобным инструментом пакетной
обработки видеоархивов, где ручная настройка каждого ролика
нецелесообразна~\cite{avisynth}. VirtualDub2 в этой связке берёт на себя
роль окна предпросмотра и фронтенда финального кодирования, наследуя
проверенный набор режимов обработки оригинального
VirtualDub~\cite{virtualdub}. Таким образом, пара «код + легковесный GUI»
обеспечивает контроль над качеством и эффективность базового монтажа.
